\documentclass[a4paper,12pt]{article}
\usepackage[utf8]{inputenc}
\usepackage{amsmath, amssymb, amsthm}

\title{Analisi 1.b - Teoria Completa}
\date{}
\begin{document}
\maketitle
\vspace{-2.5cm}
\renewcommand{\contentsname}{Indice}
\tableofcontents

% ============================
% CAPITOLO 1
% ============================
\section{Integrale di Riemann in $\mathbb{R}$}

\subsection{Costruzione}
L'integrale di Riemann nasce dall'esigenza di dare un significato rigoroso al concetto di area sotto il grafico di una funzione.  
L’idea intuitiva è quella di approssimare l’area con rettangoli: si suddivide l'intervallo $[a,b]$ in sottointervalli e si costruiscono rettangoli la cui altezza è determinata dai valori della funzione.  
Si definiscono le somme superiori e inferiori, che rappresentano due stime dell’area: la prima tende a sovrastimare, la seconda a sottostimare.  
Se i limiti coincidono, la funzione è integrabile e il valore comune è l’integrale.
Nella costruzione dell'integrale facciamo riferimento alla somma di Riemann ovvero:
\[
S(f, P, {t_i}) = \sum_{i=0}^{n-1} f(t_i) (x_{i+1} - x_i)\
\]

\subsection{Funzione integrale e Teorema fondamentale}
Data $f$ integrabile su $[a,b]$, si definisce:


\[
F(x)=\int_a^x f(t)\,dt.
\]


Se $f$ è continua, $F$ è derivabile e $F'(x)=f(x)$.  
Questo risultato collega derivazione e integrazione: l’integrale accumula i valori della funzione, la derivata restituisce istantaneamente la funzione originale. 
Il primo teorema fondamentale afferma che $F$ è continua su $[a,b]$, derivabile su $]a,b[$ e che $F'(x) = f(x)$
Dimostrazione:
\[
\frac{F(x+h)-F(x)}{h} = \frac{1}{h}\int_x^{x+h}f(t) \,dt.
\]

\[\int_x^{x+h} f(t)\,dt = f(c)h\]


\[\frac{F(x+h)-F(x)}{h} = f(c)\]
Si ha quindi che :

\[\lim_{h \to 0} f(c) = f(x)\]


Il \textbf{secondo teorema} afferma che se $F$ è una primitiva di $f$ allora:

\[\int_a^b f(x)\,dx = F(b)- F(a)\]

\subsection{Primitive}
Una primitiva di $f$ è una funzione $F$ tale che $F'(x)=f(x)$.  
Le primitive differiscono per una costante: se $F$ è primitiva, anche $F+C$ lo è.  
Questo riflette la natura della derivata, che “perde” l'informazione costante.  

\subsection{Integrazione per parti e per sostituzione}
L'integrazione per parti e per sostituzione sono due tecniche fondamentali per calcolare integrali più complessi, basate rispettivamente sulla regola del prodotto e sulla composizione di funzioni.
La formula per parti deriva dalla regola del prodotto:


\[
\int_a^b u'v = [uv]_a^b - \int_a^b uv'.
\]

L'integrazione per sostituzione serve a integrare tramite cambio di variabile. Se $f$ è continua e $\phi$ è una funzione derivabile con immagine contenuta nel dominio di $f$, allora :

\[\int f(\phi(x))\phi '(x)\, dx = \int f(u)\, du\ \text { con } \text{ }  u = \phi(x) \]


\subsection{Integrali impropri}
Gli integrali impropri si presentano quando l’intervallo è infinito o la funzione ha singolarità.  
Si definiscono tramite limiti, ad esempio:


\[
\int_a^{+\infty} f(x)\,dx = \lim_{b\to+\infty}\int_a^b f(x)\,dx.
\]


\subsection{Criteri di convergenza: confronto, convergenza assoluta, monotonia.}
Gli integrali in senso generalizzato estendono il concetto di integrale a funzioni non limitate o su intervalli infiniti. 
La loro convergenza si verifica tramite criteri specifici, come il confronto e la convergenza assoluta.

\textbf{Integrale generalizzato: }quando l'intervallo è infinito oppure quando la funzione ha discontinuità in un punto dell'intervallo 
\[\int_0^{+\infty} f(x)\,dx = \lim_{R \to +\infty}\int_a^Rf(x)\,dx\]

\textbf{Convergenza}
    \begin{itemize}
        \item \textit{Criterio di confronto:} se $0 \leq f(x) \leq g(x)$ per ogni $x$ e $\int g(x)\,dx$ converge, allora $\int f(x)\, dx$ converge
        \item \textit{Convergenza assoluta: } se $\int |f(x)| \,dx$ converge , allora anche $\int f(x)\,dx$ converge
        \item \textit{Confronto asintotico: } se $f(x) \sim g(x)$ per $x \to \infty$, cioè il loro rapporto tende a $1$, allora $f$ e $g$ hanno lo stesso comportamento integrale  
    \end{itemize}

% ============================
% CAPITOLO 2
% ============================
\section{Derivate e differenziabilità}

\subsection{Derivata direzionale}
La derivata direzionale misura la variazione di una funzione a più variabili lungo una direzione specifica, mentre le derivate parziali ne descrivono le variazioni lungo gli assi coordinati.
Descriviamo la \textbf{derivata direzionale} come :

\[
D_v f(x_0)=\lim_{t\to 0}\frac{f(x_0+tv)-f(x_0)}{t}.
\]

Essa rappresenta la velocità con cui $f$ cambia  se ci si muove da x nella direzione di $\Vec{ v}$ 
Geometricamente è la pendenza della curva ottenuta restringendo $f$ alla retta $x_0+tv$.

\subsection{Derivate parziali}
Le derivate parziali sono un caso particolare di derivata direzionale, in cui la direzione è uno degli assi del sistema di riferimento 
In $\mathbb{R}^2$:
\[\frac{\partial{f}}{\partial{x_i}}= \lim_{h \to 0}\frac{f(x_1, ...,x_i+h,...,x_n)-f(x)}{h}\]

\[
\frac{\partial f}{\partial x}(x_0,y_0),\quad \frac{\partial f}{\partial y}(x_0,y_0).
\]



\subsection{Differenziabilità}
$f$ è differenziabile in $x_0$ se può essere approssimata linearmente, ovvero, se il grafico ammette un piano tangente in quel punto 

Sia $f: R^n \to R$ una funzione. Si dice che $f$ è differenziabile in un punto $x_0 \in R^n $ se esiste una funzione lineare $L: R^n \to R$ tale che:

\[
\lim_{h \to 0} \frac{f(x_0 +h) -f(x_0) -L(h)}{||h||}
\]


Significato: il grafico è localmente “quasi lineare”, cioè ammette un piano tangente.

\subsection{Teorema}
Se $f$ è differenziabile, allora è continua e derivabile in ogni direzione:


\[
D_v f(x_0)=\nabla f(x_0)\cdot v.
\]
Rappresenta la variazione infinitesima della funzione in termini delle variazioni infinitesime delle variabili, pesate dalle derivate parziali 

\subsection{Differenziabile implica continua e derivabile in ogni direzione}
Se una funzione $f: R \to R^n$ è differenziabile in un punto $x_0$, allora è anche continua in $x_0$ e derivabile in ogni direzione in $x_0$. Questo teorema mostra che la differenziabilità è una proprietà più forte della semplice derivata direzionale
\textbf{Dimostrazione: }

\[\lim_{h \to 0} \frac{f(x_0 +h) +f(x_0) -L(h)}{||h||} = 0\]
\[f(x_0 +h) = f(x_0)+L(h)+o(||h||) \]

\subsection{Hessiana e Jacobiana}
La \textbf{Hessiana} è la matrice delle derivate seconde di $f:\mathbb{R}^n\to\mathbb{R}$, descrive la curvatura e serve a classificare i punti critici.  
La Jacobiana è la matrice delle derivate parziali di $F:\mathbb{R}^n\to\mathbb{R}^m$, rappresenta la migliore approssimazione lineare di $F$ vicino a un punto.

% ============================
% CAPITOLO 3
% ============================
\section{Integrali multipli e cambi di variabile}

\subsection{Integrale su un rettangolo}
Si costruisce come nel caso unidimensionale, suddividendo il rettangolo in blocchi elementari e considerando supremi e infimi.

La formula generale è: 
\[S(f, P) = \sum f(x_i) vol(R_i)\]  \text{dove $vol(R_i)$ è il volume dei rettangoli di $R_i$}

Se al raffinarsi, queste somme convergono a un valore indipendente dalla scelta di $x_i$, allora $f$ è integrabile secondo Riemann su $f$

\[\int_R f(x)\,dx\]
\subsection{Integrali doppi}
Se $D$ è semplice rispetto a $x$:


\[
\iint_D f(x,y)\,dx\,dy=\int_a^b\left(\int_{\varphi_1(x)}^{\varphi_2(x)} f(x,y)\,dy\right)dx.
\]

Dove $D<R^2$ è un dominio semplice in x. In questa formula si integra prima lungo una direzione tenendo l'altra fissa, e poi si somma lungo l'altra direzione.


\subsection{Integrali tripli}
Per calcolare gli integrali tripli su domini $R^3$, si usano le formule per fili e strati, che corrispondono a domini semplici rispetto a una variabile.\\ 
\textbf{Per Strati:}

\[
\iiint_D f(x,y,z)\,dx\,dy\,dz=\iint_B\left(\int_{\varphi_1(x,y)}^{\varphi_2(x,y)} f(x,y,z)\,dz\right)dx\,dy.
\]\\

\textbf{Per Fili:}

\[
\iiint_Df(x,y,z)\,dx\,dy\,dz = \int_c^d(\int_e^f(\int_{\alpha(y,z)}^{\beta(y,z)} f(x,y,z)\,dx)\,dy)\,dz
\]
Questa formula viene utilizzata quando il dominio è semplice rispetto a x, cioè per ogni $(y,z)$, la variabile $x$ varia tra due funzioni $\alpha(y,z)$ e $\beta(y,z)$
\subsection{Coordinate polari}
Le coordinate polari sono una trasformazione che esprime un punto $(x,y) \in R^2$ in termini di raggio $r \geq 0$ e angolo $\theta \in [0, 2\pi)$, utilizzando le seguenti formule:

\[
x=r\cos\theta,\ y=r\sin\theta,\quad J=r.
\]

\textbf{Jacobiano della trasformazione: }

\[
J(R, \theta) = 
\begin{bmatrix}
\frac{\partial x}{\partial r} & \frac{\partial x}{\partial \theta}\\
\frac{\partial y}{\partial r} & \frac{\partial y}{\partial \theta}
\end{bmatrix} = 
\begin{bmatrix}
cos\theta & -rsin\theta\\
sin\theta & rcos\theta
\end{bmatrix}
\]

Il valore assoluto del Jacobiano è $|J| = r$, che rappresenta il fattore di scala dell'area nel passaggio da coordinate cartesiane a polari 

Formula:


\[
\iint_D f(x,y)\,dx\,dy=\iint_{D'} f(r\cos\theta,r\sin\theta)\,r\,dr\,d\theta.
\]



\subsection{Coordinate cilindriche}
Le coordinate cilindriche sono una trasformazione utile per integrare su domini tridimensionali con simmetria circolare attorno all'asse $z$\\

\[
x=r\cos\theta,\ y=r\sin\theta,\ z=z,\quad J=r.
\]

\textbf{Jacobiano della trasformazione: }
\[
J(R, \theta, z) = 
\begin{bmatrix}
\frac{\partial x}{\partial r} & \frac{\partial x}{\partial \theta} & \frac{\partial x}{\partial z}\\
\frac{\partial y}{\partial r} & \frac{\partial y}{\partial \theta} & \frac{\partial y}{\partial z}\\
\frac{\partial z}{\partial r} & \frac{\partial z}{\partial \theta} & \frac{\partial z}{\partial z}
\end{bmatrix} = 
\begin{bmatrix}
cos\theta & -rsin\theta & 0\\
sin\theta & rcos\theta & 0 \\
0 & 0 & 1
\end{bmatrix}
\]\\
Formula cambio variabile negli integrali tripli :

\[
\iiint_D f(x,y,z)\,dx\,dy\,dz=\iiint_{D'} f(r\cos\theta,r\sin\theta,z)\,r\,dr\,d\theta\,dz.
\]

\textit{D' = dominio in coordinate cilindriche}

\subsection{Coordinate sferiche}

Le coordinate sferiche sono una trasformazione utile per integrare su domini con simmetria sferica.\\
Le formule utilizzate sono: 
\[
x=\rho\sin\phi\cos\theta,\ y=\rho\sin\phi\sin\theta,\ z=\rho\cos\phi,\quad J=\rho^2\sin\phi.
\]
\textit{$\rho$ = distanza dall'origine}\\
\textit{$\phi$ = angolo rispetto all'asse $z$}

$|J| = \rho^2sin\phi$
Formula integrale :
\[
\iiint_D f(x,y,z)\,dx\,dy\,dz=\iiint_{D'} f(\rho\sin\phi\cos\theta,\rho\sin\phi\sin\theta,\rho\cos\phi)\,\rho^2\sin\phi\,d\rho\,d\phi\,d\theta.
\]
\textit{$D'$ = dominio in coordinate sferiche}


% ============================
% CAPITOLO 4
% ============================
\section{Numeri complessi}

\subsection{Definizione}
Un numero complesso è $z=x+iy$, con $i^2=-1$.  
La parte reale è $x$, la parte immaginaria è $y$.\\
Questa tipologia di numeri consente di trattare soluzioni di equazioni come $x^2+1=0$

\subsection{Forme}
- \textbf{  Algebrica}: $z=x+iy$. \\ 
- \textbf{Trigonometrica}: $z=|z|(\cos\theta+i\sin\theta)$.\\
- \textbf{ Esponenziale}: (ricavata dalla formula di Eulero) $z=|z|e^{i\theta}$.

\subsection{Radici $n$-me}
Le radici $n$-me di $z$ sono:


\[
w_k=|z|^{1/n}\left(\cos\frac{\theta+2k\pi}{n}+i\sin\frac{\theta+2k\pi}{n}\right),\quad k=0,\dots,n-1.
\]


Dimostrazione: si scrive $w=\rho(\cos\varphi+i\sin\varphi)$, si impone $w^n=z$ e si ricavano $\rho=|z|^{1/n}$ e $\varphi=\frac{\theta+2k\pi}{n}$.  
Significato geometrico: le radici sono disposte su una circonferenza di raggio $|z|^{1/n}$, equidistanti di $2\pi/n$.

% ============================
% CAPITOLO 5
% ============================
\section{Equazioni differenziali ordinarie}

\subsection{Teorema di Cauchy di esistenza ed unicità}
Il teorema di Cauchy è uno dei risultati fondamentali della teoria delle equazioni differenziali.  
Consideriamo il problema di Cauchy:


\[
y'(t) = f(t,y(t)), \quad y(t_0) = y_0.
\]


Se $f$ è continua in un intorno di $(t_0,y_0)$ e soddisfa una condizione di Lipschitz rispetto alla variabile $y$, allora esiste un intervallo $I$ contenente $t_0$ e una funzione $y:I\to\mathbb{R}$ che risolve l’equazione. Inoltre, tale soluzione è unica.  

\textbf{Significato}: il teorema garantisce che i dati iniziali determinano una sola traiettoria. Questo è cruciale perché assicura che il modello matematico sia ben posto e non porti a soluzioni ambigue.

\subsection{Equazioni a variabili separabili}
Un'equazione differenziale si dice a variabili separabili se può essere scritta come:


\[
y'(t) = g(t)h(y).
\]


Il metodo di risoluzione consiste nel separare le variabili e integrare:


\[
\frac{1}{h(y)}\,dy = g(t)\,dt.
\]


Integrando entrambi i membri si ottiene una relazione tra $y$ e $t$.  

\textbf{Esempio}: consideriamo $y'(t) = t y$.  
Separando: $\frac{1}{y}\,dy = t\,dt$.  
Integrando: $\ln|y| = \frac{t^2}{2} + C$.  
Quindi:


\[
y(t) = Ce^{t^2/2}.
\]



\subsection{Equazioni lineari del primo ordine}
Un’equazione lineare del primo ordine ha la forma:


\[
y'(t) + a(t)y(t) = b(t),
\]


dove $a(t)$ e $b(t)$ sono funzioni continue su un intervallo $I$.  

Il metodo di risoluzione utilizza il \textbf{fattore integrante}:


\[
\mu(t) = e^{\int a(t)\,dt}.
\]


Moltiplicando l'equazione per $\mu(t)$ si ottiene:


\[
\frac{d}{dt}(\mu(t)y(t)) = \mu(t)b(t).
\]


Integrando:


\[
\mu(t)y(t) = \int \mu(t)b(t)\,dt + C.
\]


Infine:


\[
y(t) = \frac{1}{\mu(t)}\left(\int \mu(t)b(t)\,dt + C\right).
\]



\textbf{Esempio}: consideriamo $y'(t) + y(t) = e^t$.  
Qui $a(t)=1$, $b(t)=e^t$.  
Il fattore integrante è:


\[
\mu(t) = e^{\int 1\,dt} = e^t.
\]


Moltiplicando:


\[
\frac{d}{dt}(e^t y(t)) = e^{2t}.
\]


Integrando:


\[
e^t y(t) = \frac{1}{2}e^{2t} + C.
\]


Quindi:


\[
y(t) = \frac{1}{2}e^t + Ce^{-t}.
\]



\end{document}
